\section{Relation to admissibility and STOP}
\label{sec:relation-to-stop}

This section formalizes the relationship between admissibility,
falsifiability, and the \emph{STOP} outcome.
Although these notions interact operationally within diagnostic
procedures, they occupy strictly distinct logical roles in the Ontology of
Continua (OC Core) and the executable specification ETS.K\_EA.
Failure to distinguish these roles leads to category errors in the
interpretation of diagnostic results.

\medskip

\begin{definition}[Admissibility domain]
The \emph{admissibility domain} is the set of all observations, inferences,
and state transitions that are expressible without violating the axioms,
invariants, thresholds, and structural constraints defining the state space
\(\Omega(K_{EA})\).
Only elements belonging to this domain are considered meaningful inputs to
the model.
Elements outside the admissibility domain are excluded prior to any
diagnostic or falsifiability assessment.
\end{definition}

\medskip

\begin{remark}[STOP as a procedural outcome]
A \emph{STOP} outcome is a formally defined terminal state of a diagnostic
process.
It occurs when no admissible continuation of diagnosis, state refinement,
or inference exists without violating at least one mandatory model
constraint, while the existence of admissible model realizations in
\(\Omega(K_{EA})\) is not negated.
STOP is therefore a property of the diagnostic procedure rather than a
statement about the emptiness, inconsistency, or non-existence of the model
state space itself.
\end{remark}

\medskip

\begin{proposition}[Admissibility precedes falsifiability]
Falsifiability is evaluated only after admissibility has been established.
Observations that are ill-formed, externally defined, or not representable
within the formal vocabulary of ETS.K\_EA are excluded prior to any
falsifiability assessment.
Inadmissible inputs are rejected without semantic interpretation and do not
give rise to falsifying statements.
\end{proposition}

\medskip

\begin{proposition}[STOP implies non-falsification]
If a diagnostic process terminates in a STOP outcome, the
enterprise-as-continuum model is not falsified.
STOP signifies exhaustion of admissible inference paths under the given
observation interface and constraint set, not violation of model axioms,
invariants, thresholds, or structural requirements.
\end{proposition}

\medskip

\begin{proposition}[Falsification excludes STOP]
Internal falsification and STOP are mutually exclusive outcomes.
Falsification requires demonstration that no admissible model realization
exists in \(\Omega(K_{EA})\).
By contrast, STOP presupposes the continued existence of at least one
admissible realization whose internal structure cannot be further refined
without exceeding admissibility bounds.
\end{proposition}

\medskip

\begin{proposition}[Logical ordering of diagnostic outcomes]
For any diagnostic process operating within the admissibility domain, the
possible terminal outcomes are strictly ordered as follows:
\[
\text{inadmissible input}
\;\rightarrow\;
\text{admissible continuation}
\;\rightarrow\;
\begin{cases}
\text{falsification},\\
\text{STOP}.
\end{cases}
\]
No admissible diagnostic process can yield both falsification and STOP
simultaneously.
\end{proposition}

\medskip

\begin{remark}
Admissibility defines the space of meaningful interaction with the model.
Within this space, falsification corresponds to a strict logical
contradiction, namely the emptiness of the set of compatible model
realizations.
STOP, by contrast, denotes a controlled termination of inference under
constraint preservation.
STOP is therefore a first-class diagnostic outcome, not a failure mode and
not a weakened or partial form of falsification.
\end{remark}
